% Part:sets-functions-relations
% Chapter: sets
% Section: schroder-bernstein

\documentclass[../../../include/open-logic-section]{subfiles}

\begin{document}

\olfileid[fa-IR]{sfr}{siz}{sb}

\olsection{مفهوم اندازه و قضیهٔ شرودر–برنشتاین}

\begin{explain}
اندیشهٔ شهودیِ زیر را در نظر بگیرید: اگر $A$ از $B$ بزرگ‌تر نباشد و
$B$ از $A$ بزرگ‌تر نباشد، آنگاه $A$ و $B$ هم‌توان‌اند. انصافاً اگر این
اندیشه \emph{نادرست} بود، به‌دشواری می‌توانستیم این تصور را توجیه کنیم
که مفهوم تعریف‌شدهٔ ما از هم‌توانی ربطی به مقایسهٔ «اندازه»های
مجموعه‌ها دارد!{} اما خوشبختانه این اندیشهٔ شهودی درست است. قضیهٔ
شرودر–برنشتاین این امر را توجیه می‌کند.
\end{explain}

\begin{thm}[شرودر–برنشتاین]
	\ollabel{thm:schroder-bernstein}
	اگر $\cardle{A}{B}$ و $\cardle{B}{A}$،
	آنگاه $\cardeq{A}{B}$.
\end{thm}

\begin{explain}
به عبارت دیگر، اگر !!a{injection} از $A$ به~$B$ و !!a{injection} از
$B$ به~$A$ وجود داشته باشد، آنگاه !!a{bijection} از $A$ به~$B$ وجود
دارد.

بااین‌حال اثبات این نتیجه واقعاً نسبتاً \emph{دشوار} است. در واقع،
هرچند کانتور این نتیجه را بیان کرد، دیگران آن را اثبات
کردند.\footnote{برای آگاهی بیشتر از تاریخچه، برای نمونه بنگرید به
\citet[صص.~165--6]{Potter2004}.}
\oliflabeldef{sfr:cardinals:card-sb:sec}{تنها زمانی می‌توانیم قضیهٔ
شرودر–برنشتاین را \emph{اثبات} کنیم که به
\olref[sfr][cardinals][card-sb]{sec} برسیم.}{}%
فعلاً می‌توانید (و باید) آن را بی‌اثبات بپذیرید.

خوشبختانه قضیهٔ شرودر–برنشتاین \emph{درست} است و تصور ما را تأیید
می‌کند که رابطه‌هایی که تعریف کردیم، یعنی $\cardeq{A}{B}$ و
$\cardle{A}{B}$، واقعاً به «اندازه» ارتباط دارند. افزون‌براین، قضیهٔ
شرودر–برنشتاین \emph{بسیار سودمند} است. یافتنِ !!a{bijection} میان دو
مجموعهٔ هم‌توان می‌تواند دشوار باشد. قضیهٔ شرودر–برنشتاین به ما اجازه
می‌دهد مقایسه را به دو حالت تجزیه کنیم، به‌گونه‌ای که فقط لازم باشد
!!a{injection} از مجموعهٔ نخست به مجموعهٔ دوم و برعکس بیابیم.
\end{explain}
\end{document}
